% ============================================================================
%  master_table_zh.tex  —  QCD 因子化各结论严格性总表（中文版）
%  来源：research/survey/SURVEY.md（2026-10-09 文献核查）
%  依赖：common/preamble_fact.tex 中的 \tagR \tagA \tagP \tagC \tagV、booktabs、tabularx；
%        文献键见 research/survey/survey_refs.bib（INSPIRE texkey）。
%  用法：% ============================================================================
%  master_table_zh.tex  —  QCD 因子化各结论严格性总表（中文版）
%  来源：research/survey/SURVEY.md（2026-10-09 文献核查）
%  依赖：common/preamble_fact.tex 中的 \tagR \tagA \tagP \tagC \tagV、booktabs、tabularx；
%        文献键见 research/survey/survey_refs.bib（INSPIRE texkey）。
%  用法：% ============================================================================
%  master_table_zh.tex  —  QCD 因子化各结论严格性总表（中文版）
%  来源：research/survey/SURVEY.md（2026-10-09 文献核查）
%  依赖：common/preamble_fact.tex 中的 \tagR \tagA \tagP \tagC \tagV、booktabs、tabularx；
%        文献键见 research/survey/survey_refs.bib（INSPIRE texkey）。
%  用法：% ============================================================================
%  master_table_zh.tex  —  QCD 因子化各结论严格性总表（中文版）
%  来源：research/survey/SURVEY.md（2026-10-09 文献核查）
%  依赖：common/preamble_fact.tex 中的 \tagR \tagA \tagP \tagC \tagV、booktabs、tabularx；
%        文献键见 research/survey/survey_refs.bib（INSPIRE texkey）。
%  用法：\input{master_table_zh}（四个浮动表，分别覆盖条目 1–6、7–10、11–20、21–30）
% ============================================================================
\newcolumntype{Y}{>{\raggedright\arraybackslash}X}
\newcolumntype{L}[1]{>{\raggedright\arraybackslash}p{#1}}

\begin{table}[p]
\centering\footnotesize
\renewcommand{\arraystretch}{1.18}
\caption{QCD 因子化相关结论的严格性总表（一）：基础与“严格核心”（条目 1--6）。标签含义见正文；
“关键文献”栏只列最关键的已核实来源。}
\label{tab:rigor-master-1}
\begin{tabularx}{\textwidth}{@{}L{0.45cm}L{2.0cm}YL{2.6cm}L{2.35cm}Y@{}}
\toprule
\# & 主题 & 已确立的结论（要点） & 关键文献 & 标签 & 主要缺口／反例 \\
\midrule
1 & Wilson OPE、光锥 OPE &
重整化微扰论中的短距 OPE（$\phi^4$：渐近且逐圈收敛）；纯 Yang--Mills 中 OPE 在规范不变算符上闭合；电流乘积的光锥展开 &
\cite{Zimmermann:1972tv,Hollands:2011gf,Hollands:2007zg,Joglekar:1976pe,Collins:1984xc,Anikin:1978tj} &
\tagA\ （$\phi^4$、纯 YM：\tagR；非微扰：\tagC） &
含夸克 QCD：仅断言“与其他理论同样证明”（Collins 1984 §12.8）；IJJSZ 2018：“widely assumed” \\
2 & DIS 矩与 DGLAP &
$M_N=\sum_i C_{N,i}A_{N,i}+O(\Lambda^2/Q^2)$；算符定义 $\MSbar$ PDF 的 DGLAP；$\gamma_{ij}(N)=\int_0^1\!dz\,z^{N-1}P_{ij}$；核函数已知至三圈 &
\cite{Gross:1973zrg,Georgi:1974wnj,Altarelli:1977zs,Collins:1981uw,Moch:2004pa} &
\tagR\ （矩：以条目 1 为前提） &
逐矩的 $O(1/Q^2)$ 余项不给出对 $x$ 一致的控制（$x\to1$） \\
3 & 红外有限性 &
$R_{e^+e^-}$ 逐阶有限；无质量理论欧氏非例外动量格林函数红外有限；IR 安全喷注截面有限 &
\cite{Poggio:1976qr,Lowenstein:1975rg,Breitenlohner:1975hg,Sterman:1978bi,Sterman:1978bj,Sterman:1979uw} &
\tagR\ （总截面、欧氏）；\tagA\ （喷注） &
有色初态 Bloch--Nordsieck：\tagV（幂次压低）\cite{Doria:1980ak,DiLieto:1980nkq}；KLN 需初末态求和及前向散射 \cite{Frye:2018xjj} \\
4 & Landau 方程、Coleman--Norton &
夹点奇异的必要（物理区亦充分）条件；经典传播解释 &
\cite{Landau:1959fi,Coleman:1965xm,Collins:2020euz,Muhlbauer:2020kut} &
\tagA &
教科书证明不完整（Collins 2011 §5.4.3）；无质量情形的完整证明仅见预印本 \\
5 & 幂次计数、主导区域 &
幂律 $Q^p(\lambda/Q)^{\alpha(R)}$\allowbreak$\times(\lambda_S/Q)^{\beta(R)}$；主导区域目录；Glauber 夹点判据 &
\cite{Sterman:1978bi,Libby:1978qf,Libby:1978bx,Collins:2011zzd,Gardi:2022khw} &
\tagA &
Feynman 规范下主导区域不完整（附加不连通硬散射，对图求和后抵消 \cite{Labastida:1984gy}）；重叠区域需减除 \\
6 & DIS 共线因子化（图方法） &
$x$ 空间 $F=\sum_i C_i\otimes f_i+O((\Lambda/Q)^a)$，PDF 为规范不变算符矩阵元 &
\cite{Ellis:1978ty,Amati:1978by,Mueller:1978xu,Collins:1989gx,Collins:2011zzd,Collins:1998rz} &
\tagA &
非局域算符的 BRST 分类未证（Collins 2011 §11.3.2 明确假定）；K 胶子余项 \cite{Collins:2008sg} \\
\bottomrule
\end{tabularx}
\end{table}

\begin{table}[p]
\centering\footnotesize
\renewcommand{\arraystretch}{1.18}
\caption{QCD 因子化相关结论的严格性总表（二）：PDF 定义、强子--强子过程与 TMD（条目 7--10）。}
\label{tab:rigor-master-2}
\begin{tabularx}{\textwidth}{@{}L{0.45cm}L{2.0cm}YL{2.6cm}L{2.35cm}Y@{}}
\toprule
\# & 主题 & 已确立的结论（要点） & 关键文献 & 标签 & 主要缺口／反例 \\
\midrule
7 & 算符定义 PDF/FF；支撑；正定性 &
规范不变定义与重整化；$|\xi|\le1$，$f_q(-\xi)=-f_{\bar q}(\xi)$；裸密度的数密度解释 &
\cite{Collins:1981uw,Jaffe:1983hp,Collins:2011zzd} &
\tagR\ （定义、重整化、支撑）；\tagC\ （$\MSbar$ 正定性） &
正定性之争：\cite{Candido:2020yat} 对 \cite{Collins:2021vke}，及 \cite{Candido:2023ujx}（约 $Q^2\gtrsim5\,$GeV$^2$ 以上，半定量） \\
8 & Drell--Yan 与 Glauber 抵消 &
包含性 DY 首阶幂次因子化；旁观者 Glauber 交换经对切割求和与围道变形后抵消 &
\cite{Bodwin:1984hc,Collins:1985ue,Collins:1988ig,Collins:2011zzd,Diehl:2015bca} &
\tagA &
Collins 2011 §14.4.8“遗留问题”；SCET 完整证明未见发表 \cite{Rothstein:2016bsq}；喷注否决/$E_T$ 类观测量 \tagV\ \cite{Gaunt:2014ska,Zeng:2015iba}；$m^4/s^2$ 阶失效 \cite{Basu:1984ba} \\
9 & 强子对撞大 $p_T$ 单举强子/喷注、重夸克 &
唯象中普遍假定的共线因子化（含碎裂函数或喷注函数） &
\cite{Nayak:2005rt,Aybat:2008ct,Collins:1985gm} &
\tagC &
有色末态的 Glauber 抵消未证 \cite{Diehl:2015bca}；Collins 2011 §14.7；N3LO 质量奇异问题 \cite{Catani:2011st} \\
10 & TMD 因子化（DY、SIDIS、$e^+e^-$） &
$b$ 空间因子化；Collins--Soper 与 RG 演化；小 $b$ OPE 至共线 PDF &
\cite{Collins:1981uk,Collins:1984kg,Ji:2004wu,Becher:2010tm,Echevarria:2011epo,Chiu:2012ir,Collins:2017oxh} &
\tagA &
DY 依赖条目 8 的 Glauber 论证；SCET 推导在论证后舍弃 Glauber 拉氏量 \cite{Boussarie:2023izj} \\
\bottomrule
\end{tabularx}
\end{table}

\begin{table}[p]
\centering\footnotesize
\renewcommand{\arraystretch}{1.18}
\caption{QCD 因子化相关结论的严格性总表（三）：过程依赖、破坏与推广（条目 11--20）。}
\label{tab:rigor-master-3}
\begin{tabularx}{\textwidth}{@{}L{0.45cm}L{2.0cm}YL{2.6cm}L{2.35cm}Y@{}}
\toprule
\# & 主题 & 已确立的结论（要点） & 关键文献 & 标签 & 主要缺口／反例 \\
\midrule
11 & Sivers 函数变号 &
$f_{1T}^{\perp}\big|_{\rm DY}=-f_{1T}^{\perp}\big|_{\rm SIDIS}$（Boer--Mulders 同） &
\cite{Collins:2002kn,Brodsky:2002cx,Ji:2002aa,Belitsky:2002sm} &
\tagR\ （以 TMD 因子化为前提） &
实验一致但非决定性：\cite{COMPASS:2023vqt,STAR:2015vmv}；全局拟合尚不能排除同号 \cite{Bury:2020vhj} \\
12 & TMD 因子化破坏 &
背对背强子产生中标准与“广义”TMD 因子化均失效 &
\cite{Collins:2007nk,Collins:2007jp,Rogers:2010dm} &
\tagV &
模型中的显式反例（足以否定一般定理） \\
13 & 类空共线因子化破坏、超领头对数、Glauber 算符 &
类空共线极限的严格因子化在圈图层次失效；超领头对数；SCET Glauber 算符 &
\cite{Catani:2011st,Forshaw:2012bi,Forshaw:2006fk,Becher:2021zkk,Rothstein:2016bsq} &
\tagV\ （严格共线因子化）；\tagC\ （NNLO 以上质量奇异因子化） &
三圈 active--active Glauber 使间隙截面在否决标度以下与 PDF 因子化相容 \cite{Becher:2024kmk} \\
14 & 衍射因子化 &
衍射 DIS 因子化成立；强子--强子硬衍射不成立 &
\cite{Collins:1997sr,CDF:2000rua,CDF:2012nhl,H1:2007jtx} &
\tagA\ （DDIS）；\tagV\ （$p\bar p$、分解光生） &
CDF：衍射结构函数相对 HERA 预言压低 $O(10)$；H1 光生约 0.5 \\
15 & 双部分子散射 &
双 DY：全阶 Glauber 抵消（共线与 TMD）；SPS/DPS 重复计数解决；软胶子脱耦；求和规则 &
\cite{Diehl:2011yj,Diehl:2015bca,Diehl:2017kgu,Diehl:2018wfy,Diehl:2018kgr} &
\tagA\ （双 DY）；\tagC\ （有色末态） &
软矩阵 $S$ 演化与正定性无一般证明 \cite{Diehl:2017wew}；MPI 敏感观测量 \cite{Gaunt:2014ska} \\
16 & 硬遍举过程 &
$\pi$ 形状因子首阶幂次；DVCS；纵向光子介子电产生；横向光子端点失效 &
\cite{Lepage:1980fj,Efremov:1979qk,Collins:1996fb,Collins:1998be,Mankiewicz:1999tt} &
\tagA；\tagV\ （$\gamma^*_T$） &
可及 $Q^2$ 下端点与软重叠主导 \cite{Isgur:1988iw}；Sudakov/$k_T$ 方案为模型 \cite{Li:1992nu} \\
17 & SCET 因子化证明 &
算符层次软--共线脱耦；$B\to D\pi$ 全阶；DIS/DY/DVCS 重推；双喷注/thrust &
\cite{Bauer:2000yr,Bauer:2001cu,Bauer:2001yt,Bauer:2002nz,Beneke:2002ph,Manohar:2006nz} &
\tagA &
假定模式完备与 Glauber 模缺席或抵消；Collins 2011 §15.9 的保留意见 \\
18 & 无粲 $B$ 衰变 QCD 因子化（BBNS） &
首阶幂次在 $O(\alpha_s)$ 与两圈（NNLO）验证；$B\to D\pi$ 两圈证明 &
\cite{Beneke:1999br,Beneke:2000ry,Beneke:2001ev,Bell:2009nk,Bell:2020qus} &
\tagP &
非全阶证明（BBNS 2000 自述）；湮灭与扭度 3 旁观者幂次修正端点发散 \\
19 & NRQCD 因子化 &
衰变 $\Gamma=\sum_n 2\,{\rm Im}f_n$\allowbreak$\times\langle O_n\rangle/M^{d_n-4}$；产生 $\sigma=\sum_n F_n$\allowbreak$\times\langle O^H_n\rangle/M^{d_n-4}$ &
\cite{Bodwin:1994jh,Nayak:2005rw,Nayak:2005rt,Nayak:2006fm} &
\tagA\ （衰变）；\tagC\ （产生；NNLO 检验 \tagP） &
八重态产生算符需 Wilson 线；更高阶或需新非局域算符；普适性张力 \cite{Butenschoen:2012px} \\
20 & 小 $x$：$k_T$ 因子化、BFKL、BK/JIMWLK &
LL 高能因子化；LL/NLL BFKL；LL BK/JIMWLK（及 NLO BK） &
\cite{Catani:1990eg,Collins:1991ty,Fadin:1998py,Balitsky:1995ub,Kovchegov:1999yj} &
\tagP\ （$k_T$）；\tagA\ （LL/NLL 内）；\tagC\ （饱和） &
非积分胶子分布过程依赖 \cite{Dominguez:2011wm}；无受控幂次修正 \\
\bottomrule
\end{tabularx}
\end{table}

\begin{table}[p]
\centering\footnotesize
\renewcommand{\arraystretch}{1.18}
\caption{QCD 因子化相关结论的严格性总表（四）：高扭度、方法论与其他（条目 21--30）。}
\label{tab:rigor-master-4}
\begin{tabularx}{\textwidth}{@{}L{0.45cm}L{2.0cm}YL{2.6cm}L{2.35cm}Y@{}}
\toprule
\# & 主题 & 已确立的结论（要点） & 关键文献 & 标签 & 主要缺口／反例 \\
\midrule
21 & 高扭度与扭度 3 &
DIS 扭度 4（OPE）；强子 $1/Q^2$ 广义因子化；扭度 3 单自旋不对称 &
\cite{Ellis:1982cd,Jaffe:1983hp,Qiu:1990xy,Qiu:1991pp,Efremov:1984ip} &
\tagR\ （DIS 矩）；\tagA\ （$1/Q^2$）；\tagP\ （扭度 3）；\tagV\ （强子 $1/Q^4$） &
与首阶扭度的重整子混合；扭度 3 无全阶证明 \\
22 & 重整子 &
IR 重整子对应幂次修正；微扰级数阶乘发散 &
\cite{tHooft:1977xjm,Mueller:1984vh,Beneke:1998ui,Bali:2014fea} &
\tagA；\tagC\ （非微扰 OPE） &
QCD 中无非微扰的 OPE 定义 \\
23 & 区域方法 &
欧氏极限展开严格；闵氏阈值、在壳、Glauber 展开仅在特例中证明 &
\cite{Smirnov:1990rz,Beneke:1997zp,Jantzen:2011nz,Semenova:2018cwy,Gardi:2022khw} &
\tagR\ （欧氏）；\tagP\ （闵氏） &
一般情形为猜想 \cite{Semenova:2018cwy} \\
24 & LaMET 与 pseudo-PDF &
准 PDF 与 pseudo-ITD 对光锥 PDF 的匹配；Wilson 线夸克双线性算符全阶乘法可重整 &
\cite{Ji:2013dva,Izubuchi:2018srq,Radyushkin:2017cyf,Ishikawa:2017faj,Ji:2017oey} &
\tagA &
依赖 QCD 类空 OPE；幂次修正 $\sim\Lambda^2/(x^2(1-x)P_z^2)$ \cite{Braun:2018brg}；批评 \cite{Rossi:2018zkn} \\
25 & 阈值重求和与再因子化 &
首阶幂次阈值因子化与 Mellin 空间指数化；次首阶幂次 LL &
\cite{Sterman:1986aj,Catani:1989ne,Becher:2006mr,Beneke:2018gvs} &
\tagA\ （LP）；\tagP\ （NLP） &
NLP 端点发散需逐例处理 \cite{Beneke:2022obx} \\
26 & 非全局对数 &
$O(\alpha_s^2)$ 起的 NGL；BMS 方程（大 $N_c$）；有限 $N_c$；EFT 因子化 &
\cite{Dasgupta:2001sh,Banfi:2002hw,Weigert:2003mm,Becher:2016mmh} &
\tagR\ （存在性）；\tagA\ （大 $N_c$ LL）；\tagP &
强子对撞中 Glauber 相位导致超领头对数（条目 13） \\
27 & DIS 重夸克方案 &
保留重夸克质量的因子化，余项一致为 $O(\Lambda/Q)$；ACOT 等方案 &
\cite{Collins:1978wz,Aivazis:1993pi,Collins:1998rz,Kramer:2000hn,Forte:2010ta} &
\tagA &
强子--强子中保留入射夸克质量不可取（Collins 2011 §14.6） \\
28 & 光生与分解光子 &
光子结构函数（点状部分）；直接光生类 DIS；分解光生类强子--强子；质子光子 PDF &
\cite{Witten:1977ju,Bardeen:1978hg,Klasen:2002xb,Manohar:2016nzj} &
\tagA\ （直接）；\tagC\ （分解）；\tagV\ （衍射分解） &
直接/分解在 LO 以上的划分不唯一；衍射光生压低 \cite{H1:2007jtx,Klasen:2008ah} \\
29 & 核 PDF 与核中因子化 &
$eA$ 同 DIS；$pA$ DY 同 DY 加 $A^{1/3}/Q^2$ 核增强扭度 4 &
\cite{Qiu:1990xy,Luo:1994np,Bodwin:1988fs,Eskola:2021nhw} &
\tagA &
要求能量远大于随靶长度增长的标度 \cite{Bodwin:1988fs}；小 $x$ 饱和区 \tagC \\
30 & 其他要点 &
$e^+e^-\to hX$ 碎裂；Sudakov 形状因子；多腿振幅红外结构；振幅层次硬--软--共线因子化 &
\cite{Collins:1981uw,Mueller:1979ih,Collins:1980ih,Becher:2009cu,Almelid:2015jia,Feige:2014wja} &
\tagA &
偶极公式作为全阶陈述在三圈失效（\tagV）；Feige--Schwartz 排除初末态共线构型 \\
\bottomrule
\end{tabularx}
\end{table}
（四个浮动表，分别覆盖条目 1–6、7–10、11–20、21–30）
% ============================================================================
\newcolumntype{Y}{>{\raggedright\arraybackslash}X}
\newcolumntype{L}[1]{>{\raggedright\arraybackslash}p{#1}}

\begin{table}[p]
\centering\footnotesize
\renewcommand{\arraystretch}{1.18}
\caption{QCD 因子化相关结论的严格性总表（一）：基础与“严格核心”（条目 1--6）。标签含义见正文；
“关键文献”栏只列最关键的已核实来源。}
\label{tab:rigor-master-1}
\begin{tabularx}{\textwidth}{@{}L{0.45cm}L{2.0cm}YL{2.6cm}L{2.35cm}Y@{}}
\toprule
\# & 主题 & 已确立的结论（要点） & 关键文献 & 标签 & 主要缺口／反例 \\
\midrule
1 & Wilson OPE、光锥 OPE &
重整化微扰论中的短距 OPE（$\phi^4$：渐近且逐圈收敛）；纯 Yang--Mills 中 OPE 在规范不变算符上闭合；电流乘积的光锥展开 &
\cite{Zimmermann:1972tv,Hollands:2011gf,Hollands:2007zg,Joglekar:1976pe,Collins:1984xc,Anikin:1978tj} &
\tagA\ （$\phi^4$、纯 YM：\tagR；非微扰：\tagC） &
含夸克 QCD：仅断言“与其他理论同样证明”（Collins 1984 §12.8）；IJJSZ 2018：“widely assumed” \\
2 & DIS 矩与 DGLAP &
$M_N=\sum_i C_{N,i}A_{N,i}+O(\Lambda^2/Q^2)$；算符定义 $\MSbar$ PDF 的 DGLAP；$\gamma_{ij}(N)=\int_0^1\!dz\,z^{N-1}P_{ij}$；核函数已知至三圈 &
\cite{Gross:1973zrg,Georgi:1974wnj,Altarelli:1977zs,Collins:1981uw,Moch:2004pa} &
\tagR\ （矩：以条目 1 为前提） &
逐矩的 $O(1/Q^2)$ 余项不给出对 $x$ 一致的控制（$x\to1$） \\
3 & 红外有限性 &
$R_{e^+e^-}$ 逐阶有限；无质量理论欧氏非例外动量格林函数红外有限；IR 安全喷注截面有限 &
\cite{Poggio:1976qr,Lowenstein:1975rg,Breitenlohner:1975hg,Sterman:1978bi,Sterman:1978bj,Sterman:1979uw} &
\tagR\ （总截面、欧氏）；\tagA\ （喷注） &
有色初态 Bloch--Nordsieck：\tagV（幂次压低）\cite{Doria:1980ak,DiLieto:1980nkq}；KLN 需初末态求和及前向散射 \cite{Frye:2018xjj} \\
4 & Landau 方程、Coleman--Norton &
夹点奇异的必要（物理区亦充分）条件；经典传播解释 &
\cite{Landau:1959fi,Coleman:1965xm,Collins:2020euz,Muhlbauer:2020kut} &
\tagA &
教科书证明不完整（Collins 2011 §5.4.3）；无质量情形的完整证明仅见预印本 \\
5 & 幂次计数、主导区域 &
幂律 $Q^p(\lambda/Q)^{\alpha(R)}$\allowbreak$\times(\lambda_S/Q)^{\beta(R)}$；主导区域目录；Glauber 夹点判据 &
\cite{Sterman:1978bi,Libby:1978qf,Libby:1978bx,Collins:2011zzd,Gardi:2022khw} &
\tagA &
Feynman 规范下主导区域不完整（附加不连通硬散射，对图求和后抵消 \cite{Labastida:1984gy}）；重叠区域需减除 \\
6 & DIS 共线因子化（图方法） &
$x$ 空间 $F=\sum_i C_i\otimes f_i+O((\Lambda/Q)^a)$，PDF 为规范不变算符矩阵元 &
\cite{Ellis:1978ty,Amati:1978by,Mueller:1978xu,Collins:1989gx,Collins:2011zzd,Collins:1998rz} &
\tagA &
非局域算符的 BRST 分类未证（Collins 2011 §11.3.2 明确假定）；K 胶子余项 \cite{Collins:2008sg} \\
\bottomrule
\end{tabularx}
\end{table}

\begin{table}[p]
\centering\footnotesize
\renewcommand{\arraystretch}{1.18}
\caption{QCD 因子化相关结论的严格性总表（二）：PDF 定义、强子--强子过程与 TMD（条目 7--10）。}
\label{tab:rigor-master-2}
\begin{tabularx}{\textwidth}{@{}L{0.45cm}L{2.0cm}YL{2.6cm}L{2.35cm}Y@{}}
\toprule
\# & 主题 & 已确立的结论（要点） & 关键文献 & 标签 & 主要缺口／反例 \\
\midrule
7 & 算符定义 PDF/FF；支撑；正定性 &
规范不变定义与重整化；$|\xi|\le1$，$f_q(-\xi)=-f_{\bar q}(\xi)$；裸密度的数密度解释 &
\cite{Collins:1981uw,Jaffe:1983hp,Collins:2011zzd} &
\tagR\ （定义、重整化、支撑）；\tagC\ （$\MSbar$ 正定性） &
正定性之争：\cite{Candido:2020yat} 对 \cite{Collins:2021vke}，及 \cite{Candido:2023ujx}（约 $Q^2\gtrsim5\,$GeV$^2$ 以上，半定量） \\
8 & Drell--Yan 与 Glauber 抵消 &
包含性 DY 首阶幂次因子化；旁观者 Glauber 交换经对切割求和与围道变形后抵消 &
\cite{Bodwin:1984hc,Collins:1985ue,Collins:1988ig,Collins:2011zzd,Diehl:2015bca} &
\tagA &
Collins 2011 §14.4.8“遗留问题”；SCET 完整证明未见发表 \cite{Rothstein:2016bsq}；喷注否决/$E_T$ 类观测量 \tagV\ \cite{Gaunt:2014ska,Zeng:2015iba}；$m^4/s^2$ 阶失效 \cite{Basu:1984ba} \\
9 & 强子对撞大 $p_T$ 单举强子/喷注、重夸克 &
唯象中普遍假定的共线因子化（含碎裂函数或喷注函数） &
\cite{Nayak:2005rt,Aybat:2008ct,Collins:1985gm} &
\tagC &
有色末态的 Glauber 抵消未证 \cite{Diehl:2015bca}；Collins 2011 §14.7；N3LO 质量奇异问题 \cite{Catani:2011st} \\
10 & TMD 因子化（DY、SIDIS、$e^+e^-$） &
$b$ 空间因子化；Collins--Soper 与 RG 演化；小 $b$ OPE 至共线 PDF &
\cite{Collins:1981uk,Collins:1984kg,Ji:2004wu,Becher:2010tm,Echevarria:2011epo,Chiu:2012ir,Collins:2017oxh} &
\tagA &
DY 依赖条目 8 的 Glauber 论证；SCET 推导在论证后舍弃 Glauber 拉氏量 \cite{Boussarie:2023izj} \\
\bottomrule
\end{tabularx}
\end{table}

\begin{table}[p]
\centering\footnotesize
\renewcommand{\arraystretch}{1.18}
\caption{QCD 因子化相关结论的严格性总表（三）：过程依赖、破坏与推广（条目 11--20）。}
\label{tab:rigor-master-3}
\begin{tabularx}{\textwidth}{@{}L{0.45cm}L{2.0cm}YL{2.6cm}L{2.35cm}Y@{}}
\toprule
\# & 主题 & 已确立的结论（要点） & 关键文献 & 标签 & 主要缺口／反例 \\
\midrule
11 & Sivers 函数变号 &
$f_{1T}^{\perp}\big|_{\rm DY}=-f_{1T}^{\perp}\big|_{\rm SIDIS}$（Boer--Mulders 同） &
\cite{Collins:2002kn,Brodsky:2002cx,Ji:2002aa,Belitsky:2002sm} &
\tagR\ （以 TMD 因子化为前提） &
实验一致但非决定性：\cite{COMPASS:2023vqt,STAR:2015vmv}；全局拟合尚不能排除同号 \cite{Bury:2020vhj} \\
12 & TMD 因子化破坏 &
背对背强子产生中标准与“广义”TMD 因子化均失效 &
\cite{Collins:2007nk,Collins:2007jp,Rogers:2010dm} &
\tagV &
模型中的显式反例（足以否定一般定理） \\
13 & 类空共线因子化破坏、超领头对数、Glauber 算符 &
类空共线极限的严格因子化在圈图层次失效；超领头对数；SCET Glauber 算符 &
\cite{Catani:2011st,Forshaw:2012bi,Forshaw:2006fk,Becher:2021zkk,Rothstein:2016bsq} &
\tagV\ （严格共线因子化）；\tagC\ （NNLO 以上质量奇异因子化） &
三圈 active--active Glauber 使间隙截面在否决标度以下与 PDF 因子化相容 \cite{Becher:2024kmk} \\
14 & 衍射因子化 &
衍射 DIS 因子化成立；强子--强子硬衍射不成立 &
\cite{Collins:1997sr,CDF:2000rua,CDF:2012nhl,H1:2007jtx} &
\tagA\ （DDIS）；\tagV\ （$p\bar p$、分解光生） &
CDF：衍射结构函数相对 HERA 预言压低 $O(10)$；H1 光生约 0.5 \\
15 & 双部分子散射 &
双 DY：全阶 Glauber 抵消（共线与 TMD）；SPS/DPS 重复计数解决；软胶子脱耦；求和规则 &
\cite{Diehl:2011yj,Diehl:2015bca,Diehl:2017kgu,Diehl:2018wfy,Diehl:2018kgr} &
\tagA\ （双 DY）；\tagC\ （有色末态） &
软矩阵 $S$ 演化与正定性无一般证明 \cite{Diehl:2017wew}；MPI 敏感观测量 \cite{Gaunt:2014ska} \\
16 & 硬遍举过程 &
$\pi$ 形状因子首阶幂次；DVCS；纵向光子介子电产生；横向光子端点失效 &
\cite{Lepage:1980fj,Efremov:1979qk,Collins:1996fb,Collins:1998be,Mankiewicz:1999tt} &
\tagA；\tagV\ （$\gamma^*_T$） &
可及 $Q^2$ 下端点与软重叠主导 \cite{Isgur:1988iw}；Sudakov/$k_T$ 方案为模型 \cite{Li:1992nu} \\
17 & SCET 因子化证明 &
算符层次软--共线脱耦；$B\to D\pi$ 全阶；DIS/DY/DVCS 重推；双喷注/thrust &
\cite{Bauer:2000yr,Bauer:2001cu,Bauer:2001yt,Bauer:2002nz,Beneke:2002ph,Manohar:2006nz} &
\tagA &
假定模式完备与 Glauber 模缺席或抵消；Collins 2011 §15.9 的保留意见 \\
18 & 无粲 $B$ 衰变 QCD 因子化（BBNS） &
首阶幂次在 $O(\alpha_s)$ 与两圈（NNLO）验证；$B\to D\pi$ 两圈证明 &
\cite{Beneke:1999br,Beneke:2000ry,Beneke:2001ev,Bell:2009nk,Bell:2020qus} &
\tagP &
非全阶证明（BBNS 2000 自述）；湮灭与扭度 3 旁观者幂次修正端点发散 \\
19 & NRQCD 因子化 &
衰变 $\Gamma=\sum_n 2\,{\rm Im}f_n$\allowbreak$\times\langle O_n\rangle/M^{d_n-4}$；产生 $\sigma=\sum_n F_n$\allowbreak$\times\langle O^H_n\rangle/M^{d_n-4}$ &
\cite{Bodwin:1994jh,Nayak:2005rw,Nayak:2005rt,Nayak:2006fm} &
\tagA\ （衰变）；\tagC\ （产生；NNLO 检验 \tagP） &
八重态产生算符需 Wilson 线；更高阶或需新非局域算符；普适性张力 \cite{Butenschoen:2012px} \\
20 & 小 $x$：$k_T$ 因子化、BFKL、BK/JIMWLK &
LL 高能因子化；LL/NLL BFKL；LL BK/JIMWLK（及 NLO BK） &
\cite{Catani:1990eg,Collins:1991ty,Fadin:1998py,Balitsky:1995ub,Kovchegov:1999yj} &
\tagP\ （$k_T$）；\tagA\ （LL/NLL 内）；\tagC\ （饱和） &
非积分胶子分布过程依赖 \cite{Dominguez:2011wm}；无受控幂次修正 \\
\bottomrule
\end{tabularx}
\end{table}

\begin{table}[p]
\centering\footnotesize
\renewcommand{\arraystretch}{1.18}
\caption{QCD 因子化相关结论的严格性总表（四）：高扭度、方法论与其他（条目 21--30）。}
\label{tab:rigor-master-4}
\begin{tabularx}{\textwidth}{@{}L{0.45cm}L{2.0cm}YL{2.6cm}L{2.35cm}Y@{}}
\toprule
\# & 主题 & 已确立的结论（要点） & 关键文献 & 标签 & 主要缺口／反例 \\
\midrule
21 & 高扭度与扭度 3 &
DIS 扭度 4（OPE）；强子 $1/Q^2$ 广义因子化；扭度 3 单自旋不对称 &
\cite{Ellis:1982cd,Jaffe:1983hp,Qiu:1990xy,Qiu:1991pp,Efremov:1984ip} &
\tagR\ （DIS 矩）；\tagA\ （$1/Q^2$）；\tagP\ （扭度 3）；\tagV\ （强子 $1/Q^4$） &
与首阶扭度的重整子混合；扭度 3 无全阶证明 \\
22 & 重整子 &
IR 重整子对应幂次修正；微扰级数阶乘发散 &
\cite{tHooft:1977xjm,Mueller:1984vh,Beneke:1998ui,Bali:2014fea} &
\tagA；\tagC\ （非微扰 OPE） &
QCD 中无非微扰的 OPE 定义 \\
23 & 区域方法 &
欧氏极限展开严格；闵氏阈值、在壳、Glauber 展开仅在特例中证明 &
\cite{Smirnov:1990rz,Beneke:1997zp,Jantzen:2011nz,Semenova:2018cwy,Gardi:2022khw} &
\tagR\ （欧氏）；\tagP\ （闵氏） &
一般情形为猜想 \cite{Semenova:2018cwy} \\
24 & LaMET 与 pseudo-PDF &
准 PDF 与 pseudo-ITD 对光锥 PDF 的匹配；Wilson 线夸克双线性算符全阶乘法可重整 &
\cite{Ji:2013dva,Izubuchi:2018srq,Radyushkin:2017cyf,Ishikawa:2017faj,Ji:2017oey} &
\tagA &
依赖 QCD 类空 OPE；幂次修正 $\sim\Lambda^2/(x^2(1-x)P_z^2)$ \cite{Braun:2018brg}；批评 \cite{Rossi:2018zkn} \\
25 & 阈值重求和与再因子化 &
首阶幂次阈值因子化与 Mellin 空间指数化；次首阶幂次 LL &
\cite{Sterman:1986aj,Catani:1989ne,Becher:2006mr,Beneke:2018gvs} &
\tagA\ （LP）；\tagP\ （NLP） &
NLP 端点发散需逐例处理 \cite{Beneke:2022obx} \\
26 & 非全局对数 &
$O(\alpha_s^2)$ 起的 NGL；BMS 方程（大 $N_c$）；有限 $N_c$；EFT 因子化 &
\cite{Dasgupta:2001sh,Banfi:2002hw,Weigert:2003mm,Becher:2016mmh} &
\tagR\ （存在性）；\tagA\ （大 $N_c$ LL）；\tagP &
强子对撞中 Glauber 相位导致超领头对数（条目 13） \\
27 & DIS 重夸克方案 &
保留重夸克质量的因子化，余项一致为 $O(\Lambda/Q)$；ACOT 等方案 &
\cite{Collins:1978wz,Aivazis:1993pi,Collins:1998rz,Kramer:2000hn,Forte:2010ta} &
\tagA &
强子--强子中保留入射夸克质量不可取（Collins 2011 §14.6） \\
28 & 光生与分解光子 &
光子结构函数（点状部分）；直接光生类 DIS；分解光生类强子--强子；质子光子 PDF &
\cite{Witten:1977ju,Bardeen:1978hg,Klasen:2002xb,Manohar:2016nzj} &
\tagA\ （直接）；\tagC\ （分解）；\tagV\ （衍射分解） &
直接/分解在 LO 以上的划分不唯一；衍射光生压低 \cite{H1:2007jtx,Klasen:2008ah} \\
29 & 核 PDF 与核中因子化 &
$eA$ 同 DIS；$pA$ DY 同 DY 加 $A^{1/3}/Q^2$ 核增强扭度 4 &
\cite{Qiu:1990xy,Luo:1994np,Bodwin:1988fs,Eskola:2021nhw} &
\tagA &
要求能量远大于随靶长度增长的标度 \cite{Bodwin:1988fs}；小 $x$ 饱和区 \tagC \\
30 & 其他要点 &
$e^+e^-\to hX$ 碎裂；Sudakov 形状因子；多腿振幅红外结构；振幅层次硬--软--共线因子化 &
\cite{Collins:1981uw,Mueller:1979ih,Collins:1980ih,Becher:2009cu,Almelid:2015jia,Feige:2014wja} &
\tagA &
偶极公式作为全阶陈述在三圈失效（\tagV）；Feige--Schwartz 排除初末态共线构型 \\
\bottomrule
\end{tabularx}
\end{table}
（四个浮动表，分别覆盖条目 1–6、7–10、11–20、21–30）
% ============================================================================
\newcolumntype{Y}{>{\raggedright\arraybackslash}X}
\newcolumntype{L}[1]{>{\raggedright\arraybackslash}p{#1}}

\begin{table}[p]
\centering\footnotesize
\renewcommand{\arraystretch}{1.18}
\caption{QCD 因子化相关结论的严格性总表（一）：基础与“严格核心”（条目 1--6）。标签含义见正文；
“关键文献”栏只列最关键的已核实来源。}
\label{tab:rigor-master-1}
\begin{tabularx}{\textwidth}{@{}L{0.45cm}L{2.0cm}YL{2.6cm}L{2.35cm}Y@{}}
\toprule
\# & 主题 & 已确立的结论（要点） & 关键文献 & 标签 & 主要缺口／反例 \\
\midrule
1 & Wilson OPE、光锥 OPE &
重整化微扰论中的短距 OPE（$\phi^4$：渐近且逐圈收敛）；纯 Yang--Mills 中 OPE 在规范不变算符上闭合；电流乘积的光锥展开 &
\cite{Zimmermann:1972tv,Hollands:2011gf,Hollands:2007zg,Joglekar:1976pe,Collins:1984xc,Anikin:1978tj} &
\tagA\ （$\phi^4$、纯 YM：\tagR；非微扰：\tagC） &
含夸克 QCD：仅断言“与其他理论同样证明”（Collins 1984 §12.8）；IJJSZ 2018：“widely assumed” \\
2 & DIS 矩与 DGLAP &
$M_N=\sum_i C_{N,i}A_{N,i}+O(\Lambda^2/Q^2)$；算符定义 $\MSbar$ PDF 的 DGLAP；$\gamma_{ij}(N)=\int_0^1\!dz\,z^{N-1}P_{ij}$；核函数已知至三圈 &
\cite{Gross:1973zrg,Georgi:1974wnj,Altarelli:1977zs,Collins:1981uw,Moch:2004pa} &
\tagR\ （矩：以条目 1 为前提） &
逐矩的 $O(1/Q^2)$ 余项不给出对 $x$ 一致的控制（$x\to1$） \\
3 & 红外有限性 &
$R_{e^+e^-}$ 逐阶有限；无质量理论欧氏非例外动量格林函数红外有限；IR 安全喷注截面有限 &
\cite{Poggio:1976qr,Lowenstein:1975rg,Breitenlohner:1975hg,Sterman:1978bi,Sterman:1978bj,Sterman:1979uw} &
\tagR\ （总截面、欧氏）；\tagA\ （喷注） &
有色初态 Bloch--Nordsieck：\tagV（幂次压低）\cite{Doria:1980ak,DiLieto:1980nkq}；KLN 需初末态求和及前向散射 \cite{Frye:2018xjj} \\
4 & Landau 方程、Coleman--Norton &
夹点奇异的必要（物理区亦充分）条件；经典传播解释 &
\cite{Landau:1959fi,Coleman:1965xm,Collins:2020euz,Muhlbauer:2020kut} &
\tagA &
教科书证明不完整（Collins 2011 §5.4.3）；无质量情形的完整证明仅见预印本 \\
5 & 幂次计数、主导区域 &
幂律 $Q^p(\lambda/Q)^{\alpha(R)}$\allowbreak$\times(\lambda_S/Q)^{\beta(R)}$；主导区域目录；Glauber 夹点判据 &
\cite{Sterman:1978bi,Libby:1978qf,Libby:1978bx,Collins:2011zzd,Gardi:2022khw} &
\tagA &
Feynman 规范下主导区域不完整（附加不连通硬散射，对图求和后抵消 \cite{Labastida:1984gy}）；重叠区域需减除 \\
6 & DIS 共线因子化（图方法） &
$x$ 空间 $F=\sum_i C_i\otimes f_i+O((\Lambda/Q)^a)$，PDF 为规范不变算符矩阵元 &
\cite{Ellis:1978ty,Amati:1978by,Mueller:1978xu,Collins:1989gx,Collins:2011zzd,Collins:1998rz} &
\tagA &
非局域算符的 BRST 分类未证（Collins 2011 §11.3.2 明确假定）；K 胶子余项 \cite{Collins:2008sg} \\
\bottomrule
\end{tabularx}
\end{table}

\begin{table}[p]
\centering\footnotesize
\renewcommand{\arraystretch}{1.18}
\caption{QCD 因子化相关结论的严格性总表（二）：PDF 定义、强子--强子过程与 TMD（条目 7--10）。}
\label{tab:rigor-master-2}
\begin{tabularx}{\textwidth}{@{}L{0.45cm}L{2.0cm}YL{2.6cm}L{2.35cm}Y@{}}
\toprule
\# & 主题 & 已确立的结论（要点） & 关键文献 & 标签 & 主要缺口／反例 \\
\midrule
7 & 算符定义 PDF/FF；支撑；正定性 &
规范不变定义与重整化；$|\xi|\le1$，$f_q(-\xi)=-f_{\bar q}(\xi)$；裸密度的数密度解释 &
\cite{Collins:1981uw,Jaffe:1983hp,Collins:2011zzd} &
\tagR\ （定义、重整化、支撑）；\tagC\ （$\MSbar$ 正定性） &
正定性之争：\cite{Candido:2020yat} 对 \cite{Collins:2021vke}，及 \cite{Candido:2023ujx}（约 $Q^2\gtrsim5\,$GeV$^2$ 以上，半定量） \\
8 & Drell--Yan 与 Glauber 抵消 &
包含性 DY 首阶幂次因子化；旁观者 Glauber 交换经对切割求和与围道变形后抵消 &
\cite{Bodwin:1984hc,Collins:1985ue,Collins:1988ig,Collins:2011zzd,Diehl:2015bca} &
\tagA &
Collins 2011 §14.4.8“遗留问题”；SCET 完整证明未见发表 \cite{Rothstein:2016bsq}；喷注否决/$E_T$ 类观测量 \tagV\ \cite{Gaunt:2014ska,Zeng:2015iba}；$m^4/s^2$ 阶失效 \cite{Basu:1984ba} \\
9 & 强子对撞大 $p_T$ 单举强子/喷注、重夸克 &
唯象中普遍假定的共线因子化（含碎裂函数或喷注函数） &
\cite{Nayak:2005rt,Aybat:2008ct,Collins:1985gm} &
\tagC &
有色末态的 Glauber 抵消未证 \cite{Diehl:2015bca}；Collins 2011 §14.7；N3LO 质量奇异问题 \cite{Catani:2011st} \\
10 & TMD 因子化（DY、SIDIS、$e^+e^-$） &
$b$ 空间因子化；Collins--Soper 与 RG 演化；小 $b$ OPE 至共线 PDF &
\cite{Collins:1981uk,Collins:1984kg,Ji:2004wu,Becher:2010tm,Echevarria:2011epo,Chiu:2012ir,Collins:2017oxh} &
\tagA &
DY 依赖条目 8 的 Glauber 论证；SCET 推导在论证后舍弃 Glauber 拉氏量 \cite{Boussarie:2023izj} \\
\bottomrule
\end{tabularx}
\end{table}

\begin{table}[p]
\centering\footnotesize
\renewcommand{\arraystretch}{1.18}
\caption{QCD 因子化相关结论的严格性总表（三）：过程依赖、破坏与推广（条目 11--20）。}
\label{tab:rigor-master-3}
\begin{tabularx}{\textwidth}{@{}L{0.45cm}L{2.0cm}YL{2.6cm}L{2.35cm}Y@{}}
\toprule
\# & 主题 & 已确立的结论（要点） & 关键文献 & 标签 & 主要缺口／反例 \\
\midrule
11 & Sivers 函数变号 &
$f_{1T}^{\perp}\big|_{\rm DY}=-f_{1T}^{\perp}\big|_{\rm SIDIS}$（Boer--Mulders 同） &
\cite{Collins:2002kn,Brodsky:2002cx,Ji:2002aa,Belitsky:2002sm} &
\tagR\ （以 TMD 因子化为前提） &
实验一致但非决定性：\cite{COMPASS:2023vqt,STAR:2015vmv}；全局拟合尚不能排除同号 \cite{Bury:2020vhj} \\
12 & TMD 因子化破坏 &
背对背强子产生中标准与“广义”TMD 因子化均失效 &
\cite{Collins:2007nk,Collins:2007jp,Rogers:2010dm} &
\tagV &
模型中的显式反例（足以否定一般定理） \\
13 & 类空共线因子化破坏、超领头对数、Glauber 算符 &
类空共线极限的严格因子化在圈图层次失效；超领头对数；SCET Glauber 算符 &
\cite{Catani:2011st,Forshaw:2012bi,Forshaw:2006fk,Becher:2021zkk,Rothstein:2016bsq} &
\tagV\ （严格共线因子化）；\tagC\ （NNLO 以上质量奇异因子化） &
三圈 active--active Glauber 使间隙截面在否决标度以下与 PDF 因子化相容 \cite{Becher:2024kmk} \\
14 & 衍射因子化 &
衍射 DIS 因子化成立；强子--强子硬衍射不成立 &
\cite{Collins:1997sr,CDF:2000rua,CDF:2012nhl,H1:2007jtx} &
\tagA\ （DDIS）；\tagV\ （$p\bar p$、分解光生） &
CDF：衍射结构函数相对 HERA 预言压低 $O(10)$；H1 光生约 0.5 \\
15 & 双部分子散射 &
双 DY：全阶 Glauber 抵消（共线与 TMD）；SPS/DPS 重复计数解决；软胶子脱耦；求和规则 &
\cite{Diehl:2011yj,Diehl:2015bca,Diehl:2017kgu,Diehl:2018wfy,Diehl:2018kgr} &
\tagA\ （双 DY）；\tagC\ （有色末态） &
软矩阵 $S$ 演化与正定性无一般证明 \cite{Diehl:2017wew}；MPI 敏感观测量 \cite{Gaunt:2014ska} \\
16 & 硬遍举过程 &
$\pi$ 形状因子首阶幂次；DVCS；纵向光子介子电产生；横向光子端点失效 &
\cite{Lepage:1980fj,Efremov:1979qk,Collins:1996fb,Collins:1998be,Mankiewicz:1999tt} &
\tagA；\tagV\ （$\gamma^*_T$） &
可及 $Q^2$ 下端点与软重叠主导 \cite{Isgur:1988iw}；Sudakov/$k_T$ 方案为模型 \cite{Li:1992nu} \\
17 & SCET 因子化证明 &
算符层次软--共线脱耦；$B\to D\pi$ 全阶；DIS/DY/DVCS 重推；双喷注/thrust &
\cite{Bauer:2000yr,Bauer:2001cu,Bauer:2001yt,Bauer:2002nz,Beneke:2002ph,Manohar:2006nz} &
\tagA &
假定模式完备与 Glauber 模缺席或抵消；Collins 2011 §15.9 的保留意见 \\
18 & 无粲 $B$ 衰变 QCD 因子化（BBNS） &
首阶幂次在 $O(\alpha_s)$ 与两圈（NNLO）验证；$B\to D\pi$ 两圈证明 &
\cite{Beneke:1999br,Beneke:2000ry,Beneke:2001ev,Bell:2009nk,Bell:2020qus} &
\tagP &
非全阶证明（BBNS 2000 自述）；湮灭与扭度 3 旁观者幂次修正端点发散 \\
19 & NRQCD 因子化 &
衰变 $\Gamma=\sum_n 2\,{\rm Im}f_n$\allowbreak$\times\langle O_n\rangle/M^{d_n-4}$；产生 $\sigma=\sum_n F_n$\allowbreak$\times\langle O^H_n\rangle/M^{d_n-4}$ &
\cite{Bodwin:1994jh,Nayak:2005rw,Nayak:2005rt,Nayak:2006fm} &
\tagA\ （衰变）；\tagC\ （产生；NNLO 检验 \tagP） &
八重态产生算符需 Wilson 线；更高阶或需新非局域算符；普适性张力 \cite{Butenschoen:2012px} \\
20 & 小 $x$：$k_T$ 因子化、BFKL、BK/JIMWLK &
LL 高能因子化；LL/NLL BFKL；LL BK/JIMWLK（及 NLO BK） &
\cite{Catani:1990eg,Collins:1991ty,Fadin:1998py,Balitsky:1995ub,Kovchegov:1999yj} &
\tagP\ （$k_T$）；\tagA\ （LL/NLL 内）；\tagC\ （饱和） &
非积分胶子分布过程依赖 \cite{Dominguez:2011wm}；无受控幂次修正 \\
\bottomrule
\end{tabularx}
\end{table}

\begin{table}[p]
\centering\footnotesize
\renewcommand{\arraystretch}{1.18}
\caption{QCD 因子化相关结论的严格性总表（四）：高扭度、方法论与其他（条目 21--30）。}
\label{tab:rigor-master-4}
\begin{tabularx}{\textwidth}{@{}L{0.45cm}L{2.0cm}YL{2.6cm}L{2.35cm}Y@{}}
\toprule
\# & 主题 & 已确立的结论（要点） & 关键文献 & 标签 & 主要缺口／反例 \\
\midrule
21 & 高扭度与扭度 3 &
DIS 扭度 4（OPE）；强子 $1/Q^2$ 广义因子化；扭度 3 单自旋不对称 &
\cite{Ellis:1982cd,Jaffe:1983hp,Qiu:1990xy,Qiu:1991pp,Efremov:1984ip} &
\tagR\ （DIS 矩）；\tagA\ （$1/Q^2$）；\tagP\ （扭度 3）；\tagV\ （强子 $1/Q^4$） &
与首阶扭度的重整子混合；扭度 3 无全阶证明 \\
22 & 重整子 &
IR 重整子对应幂次修正；微扰级数阶乘发散 &
\cite{tHooft:1977xjm,Mueller:1984vh,Beneke:1998ui,Bali:2014fea} &
\tagA；\tagC\ （非微扰 OPE） &
QCD 中无非微扰的 OPE 定义 \\
23 & 区域方法 &
欧氏极限展开严格；闵氏阈值、在壳、Glauber 展开仅在特例中证明 &
\cite{Smirnov:1990rz,Beneke:1997zp,Jantzen:2011nz,Semenova:2018cwy,Gardi:2022khw} &
\tagR\ （欧氏）；\tagP\ （闵氏） &
一般情形为猜想 \cite{Semenova:2018cwy} \\
24 & LaMET 与 pseudo-PDF &
准 PDF 与 pseudo-ITD 对光锥 PDF 的匹配；Wilson 线夸克双线性算符全阶乘法可重整 &
\cite{Ji:2013dva,Izubuchi:2018srq,Radyushkin:2017cyf,Ishikawa:2017faj,Ji:2017oey} &
\tagA &
依赖 QCD 类空 OPE；幂次修正 $\sim\Lambda^2/(x^2(1-x)P_z^2)$ \cite{Braun:2018brg}；批评 \cite{Rossi:2018zkn} \\
25 & 阈值重求和与再因子化 &
首阶幂次阈值因子化与 Mellin 空间指数化；次首阶幂次 LL &
\cite{Sterman:1986aj,Catani:1989ne,Becher:2006mr,Beneke:2018gvs} &
\tagA\ （LP）；\tagP\ （NLP） &
NLP 端点发散需逐例处理 \cite{Beneke:2022obx} \\
26 & 非全局对数 &
$O(\alpha_s^2)$ 起的 NGL；BMS 方程（大 $N_c$）；有限 $N_c$；EFT 因子化 &
\cite{Dasgupta:2001sh,Banfi:2002hw,Weigert:2003mm,Becher:2016mmh} &
\tagR\ （存在性）；\tagA\ （大 $N_c$ LL）；\tagP &
强子对撞中 Glauber 相位导致超领头对数（条目 13） \\
27 & DIS 重夸克方案 &
保留重夸克质量的因子化，余项一致为 $O(\Lambda/Q)$；ACOT 等方案 &
\cite{Collins:1978wz,Aivazis:1993pi,Collins:1998rz,Kramer:2000hn,Forte:2010ta} &
\tagA &
强子--强子中保留入射夸克质量不可取（Collins 2011 §14.6） \\
28 & 光生与分解光子 &
光子结构函数（点状部分）；直接光生类 DIS；分解光生类强子--强子；质子光子 PDF &
\cite{Witten:1977ju,Bardeen:1978hg,Klasen:2002xb,Manohar:2016nzj} &
\tagA\ （直接）；\tagC\ （分解）；\tagV\ （衍射分解） &
直接/分解在 LO 以上的划分不唯一；衍射光生压低 \cite{H1:2007jtx,Klasen:2008ah} \\
29 & 核 PDF 与核中因子化 &
$eA$ 同 DIS；$pA$ DY 同 DY 加 $A^{1/3}/Q^2$ 核增强扭度 4 &
\cite{Qiu:1990xy,Luo:1994np,Bodwin:1988fs,Eskola:2021nhw} &
\tagA &
要求能量远大于随靶长度增长的标度 \cite{Bodwin:1988fs}；小 $x$ 饱和区 \tagC \\
30 & 其他要点 &
$e^+e^-\to hX$ 碎裂；Sudakov 形状因子；多腿振幅红外结构；振幅层次硬--软--共线因子化 &
\cite{Collins:1981uw,Mueller:1979ih,Collins:1980ih,Becher:2009cu,Almelid:2015jia,Feige:2014wja} &
\tagA &
偶极公式作为全阶陈述在三圈失效（\tagV）；Feige--Schwartz 排除初末态共线构型 \\
\bottomrule
\end{tabularx}
\end{table}
（四个浮动表，分别覆盖条目 1–6、7–10、11–20、21–30）
% ============================================================================
\newcolumntype{Y}{>{\raggedright\arraybackslash}X}
\newcolumntype{L}[1]{>{\raggedright\arraybackslash}p{#1}}

\begin{table}[p]
\centering\footnotesize
\renewcommand{\arraystretch}{1.18}
\caption{QCD 因子化相关结论的严格性总表（一）：基础与“严格核心”（条目 1--6）。标签含义见正文；
“关键文献”栏只列最关键的已核实来源。}
\label{tab:rigor-master-1}
\begin{tabularx}{\textwidth}{@{}L{0.45cm}L{2.0cm}YL{2.6cm}L{2.35cm}Y@{}}
\toprule
\# & 主题 & 已确立的结论（要点） & 关键文献 & 标签 & 主要缺口／反例 \\
\midrule
1 & Wilson OPE、光锥 OPE &
重整化微扰论中的短距 OPE（$\phi^4$：渐近且逐圈收敛）；纯 Yang--Mills 中 OPE 在规范不变算符上闭合；电流乘积的光锥展开 &
\cite{Zimmermann:1972tv,Hollands:2011gf,Hollands:2007zg,Joglekar:1976pe,Collins:1984xc,Anikin:1978tj} &
\tagA\ （$\phi^4$、纯 YM：\tagR；非微扰：\tagC） &
含夸克 QCD：仅断言“与其他理论同样证明”（Collins 1984 §12.8）；IJJSZ 2018：“widely assumed” \\
2 & DIS 矩与 DGLAP &
$M_N=\sum_i C_{N,i}A_{N,i}+O(\Lambda^2/Q^2)$；算符定义 $\MSbar$ PDF 的 DGLAP；$\gamma_{ij}(N)=\int_0^1\!dz\,z^{N-1}P_{ij}$；核函数已知至三圈 &
\cite{Gross:1973zrg,Georgi:1974wnj,Altarelli:1977zs,Collins:1981uw,Moch:2004pa} &
\tagR\ （矩：以条目 1 为前提） &
逐矩的 $O(1/Q^2)$ 余项不给出对 $x$ 一致的控制（$x\to1$） \\
3 & 红外有限性 &
$R_{e^+e^-}$ 逐阶有限；无质量理论欧氏非例外动量格林函数红外有限；IR 安全喷注截面有限 &
\cite{Poggio:1976qr,Lowenstein:1975rg,Breitenlohner:1975hg,Sterman:1978bi,Sterman:1978bj,Sterman:1979uw} &
\tagR\ （总截面、欧氏）；\tagA\ （喷注） &
有色初态 Bloch--Nordsieck：\tagV（幂次压低）\cite{Doria:1980ak,DiLieto:1980nkq}；KLN 需初末态求和及前向散射 \cite{Frye:2018xjj} \\
4 & Landau 方程、Coleman--Norton &
夹点奇异的必要（物理区亦充分）条件；经典传播解释 &
\cite{Landau:1959fi,Coleman:1965xm,Collins:2020euz,Muhlbauer:2020kut} &
\tagA &
教科书证明不完整（Collins 2011 §5.4.3）；无质量情形的完整证明仅见预印本 \\
5 & 幂次计数、主导区域 &
幂律 $Q^p(\lambda/Q)^{\alpha(R)}$\allowbreak$\times(\lambda_S/Q)^{\beta(R)}$；主导区域目录；Glauber 夹点判据 &
\cite{Sterman:1978bi,Libby:1978qf,Libby:1978bx,Collins:2011zzd,Gardi:2022khw} &
\tagA &
Feynman 规范下主导区域不完整（附加不连通硬散射，对图求和后抵消 \cite{Labastida:1984gy}）；重叠区域需减除 \\
6 & DIS 共线因子化（图方法） &
$x$ 空间 $F=\sum_i C_i\otimes f_i+O((\Lambda/Q)^a)$，PDF 为规范不变算符矩阵元 &
\cite{Ellis:1978ty,Amati:1978by,Mueller:1978xu,Collins:1989gx,Collins:2011zzd,Collins:1998rz} &
\tagA &
非局域算符的 BRST 分类未证（Collins 2011 §11.3.2 明确假定）；K 胶子余项 \cite{Collins:2008sg} \\
\bottomrule
\end{tabularx}
\end{table}

\begin{table}[p]
\centering\footnotesize
\renewcommand{\arraystretch}{1.18}
\caption{QCD 因子化相关结论的严格性总表（二）：PDF 定义、强子--强子过程与 TMD（条目 7--10）。}
\label{tab:rigor-master-2}
\begin{tabularx}{\textwidth}{@{}L{0.45cm}L{2.0cm}YL{2.6cm}L{2.35cm}Y@{}}
\toprule
\# & 主题 & 已确立的结论（要点） & 关键文献 & 标签 & 主要缺口／反例 \\
\midrule
7 & 算符定义 PDF/FF；支撑；正定性 &
规范不变定义与重整化；$|\xi|\le1$，$f_q(-\xi)=-f_{\bar q}(\xi)$；裸密度的数密度解释 &
\cite{Collins:1981uw,Jaffe:1983hp,Collins:2011zzd} &
\tagR\ （定义、重整化、支撑）；\tagC\ （$\MSbar$ 正定性） &
正定性之争：\cite{Candido:2020yat} 对 \cite{Collins:2021vke}，及 \cite{Candido:2023ujx}（约 $Q^2\gtrsim5\,$GeV$^2$ 以上，半定量） \\
8 & Drell--Yan 与 Glauber 抵消 &
包含性 DY 首阶幂次因子化；旁观者 Glauber 交换经对切割求和与围道变形后抵消 &
\cite{Bodwin:1984hc,Collins:1985ue,Collins:1988ig,Collins:2011zzd,Diehl:2015bca} &
\tagA &
Collins 2011 §14.4.8“遗留问题”；SCET 完整证明未见发表 \cite{Rothstein:2016bsq}；喷注否决/$E_T$ 类观测量 \tagV\ \cite{Gaunt:2014ska,Zeng:2015iba}；$m^4/s^2$ 阶失效 \cite{Basu:1984ba} \\
9 & 强子对撞大 $p_T$ 单举强子/喷注、重夸克 &
唯象中普遍假定的共线因子化（含碎裂函数或喷注函数） &
\cite{Nayak:2005rt,Aybat:2008ct,Collins:1985gm} &
\tagC &
有色末态的 Glauber 抵消未证 \cite{Diehl:2015bca}；Collins 2011 §14.7；N3LO 质量奇异问题 \cite{Catani:2011st} \\
10 & TMD 因子化（DY、SIDIS、$e^+e^-$） &
$b$ 空间因子化；Collins--Soper 与 RG 演化；小 $b$ OPE 至共线 PDF &
\cite{Collins:1981uk,Collins:1984kg,Ji:2004wu,Becher:2010tm,Echevarria:2011epo,Chiu:2012ir,Collins:2017oxh} &
\tagA &
DY 依赖条目 8 的 Glauber 论证；SCET 推导在论证后舍弃 Glauber 拉氏量 \cite{Boussarie:2023izj} \\
\bottomrule
\end{tabularx}
\end{table}

\begin{table}[p]
\centering\footnotesize
\renewcommand{\arraystretch}{1.18}
\caption{QCD 因子化相关结论的严格性总表（三）：过程依赖、破坏与推广（条目 11--20）。}
\label{tab:rigor-master-3}
\begin{tabularx}{\textwidth}{@{}L{0.45cm}L{2.0cm}YL{2.6cm}L{2.35cm}Y@{}}
\toprule
\# & 主题 & 已确立的结论（要点） & 关键文献 & 标签 & 主要缺口／反例 \\
\midrule
11 & Sivers 函数变号 &
$f_{1T}^{\perp}\big|_{\rm DY}=-f_{1T}^{\perp}\big|_{\rm SIDIS}$（Boer--Mulders 同） &
\cite{Collins:2002kn,Brodsky:2002cx,Ji:2002aa,Belitsky:2002sm} &
\tagR\ （以 TMD 因子化为前提） &
实验一致但非决定性：\cite{COMPASS:2023vqt,STAR:2015vmv}；全局拟合尚不能排除同号 \cite{Bury:2020vhj} \\
12 & TMD 因子化破坏 &
背对背强子产生中标准与“广义”TMD 因子化均失效 &
\cite{Collins:2007nk,Collins:2007jp,Rogers:2010dm} &
\tagV &
模型中的显式反例（足以否定一般定理） \\
13 & 类空共线因子化破坏、超领头对数、Glauber 算符 &
类空共线极限的严格因子化在圈图层次失效；超领头对数；SCET Glauber 算符 &
\cite{Catani:2011st,Forshaw:2012bi,Forshaw:2006fk,Becher:2021zkk,Rothstein:2016bsq} &
\tagV\ （严格共线因子化）；\tagC\ （NNLO 以上质量奇异因子化） &
三圈 active--active Glauber 使间隙截面在否决标度以下与 PDF 因子化相容 \cite{Becher:2024kmk} \\
14 & 衍射因子化 &
衍射 DIS 因子化成立；强子--强子硬衍射不成立 &
\cite{Collins:1997sr,CDF:2000rua,CDF:2012nhl,H1:2007jtx} &
\tagA\ （DDIS）；\tagV\ （$p\bar p$、分解光生） &
CDF：衍射结构函数相对 HERA 预言压低 $O(10)$；H1 光生约 0.5 \\
15 & 双部分子散射 &
双 DY：全阶 Glauber 抵消（共线与 TMD）；SPS/DPS 重复计数解决；软胶子脱耦；求和规则 &
\cite{Diehl:2011yj,Diehl:2015bca,Diehl:2017kgu,Diehl:2018wfy,Diehl:2018kgr} &
\tagA\ （双 DY）；\tagC\ （有色末态） &
软矩阵 $S$ 演化与正定性无一般证明 \cite{Diehl:2017wew}；MPI 敏感观测量 \cite{Gaunt:2014ska} \\
16 & 硬遍举过程 &
$\pi$ 形状因子首阶幂次；DVCS；纵向光子介子电产生；横向光子端点失效 &
\cite{Lepage:1980fj,Efremov:1979qk,Collins:1996fb,Collins:1998be,Mankiewicz:1999tt} &
\tagA；\tagV\ （$\gamma^*_T$） &
可及 $Q^2$ 下端点与软重叠主导 \cite{Isgur:1988iw}；Sudakov/$k_T$ 方案为模型 \cite{Li:1992nu} \\
17 & SCET 因子化证明 &
算符层次软--共线脱耦；$B\to D\pi$ 全阶；DIS/DY/DVCS 重推；双喷注/thrust &
\cite{Bauer:2000yr,Bauer:2001cu,Bauer:2001yt,Bauer:2002nz,Beneke:2002ph,Manohar:2006nz} &
\tagA &
假定模式完备与 Glauber 模缺席或抵消；Collins 2011 §15.9 的保留意见 \\
18 & 无粲 $B$ 衰变 QCD 因子化（BBNS） &
首阶幂次在 $O(\alpha_s)$ 与两圈（NNLO）验证；$B\to D\pi$ 两圈证明 &
\cite{Beneke:1999br,Beneke:2000ry,Beneke:2001ev,Bell:2009nk,Bell:2020qus} &
\tagP &
非全阶证明（BBNS 2000 自述）；湮灭与扭度 3 旁观者幂次修正端点发散 \\
19 & NRQCD 因子化 &
衰变 $\Gamma=\sum_n 2\,{\rm Im}f_n$\allowbreak$\times\langle O_n\rangle/M^{d_n-4}$；产生 $\sigma=\sum_n F_n$\allowbreak$\times\langle O^H_n\rangle/M^{d_n-4}$ &
\cite{Bodwin:1994jh,Nayak:2005rw,Nayak:2005rt,Nayak:2006fm} &
\tagA\ （衰变）；\tagC\ （产生；NNLO 检验 \tagP） &
八重态产生算符需 Wilson 线；更高阶或需新非局域算符；普适性张力 \cite{Butenschoen:2012px} \\
20 & 小 $x$：$k_T$ 因子化、BFKL、BK/JIMWLK &
LL 高能因子化；LL/NLL BFKL；LL BK/JIMWLK（及 NLO BK） &
\cite{Catani:1990eg,Collins:1991ty,Fadin:1998py,Balitsky:1995ub,Kovchegov:1999yj} &
\tagP\ （$k_T$）；\tagA\ （LL/NLL 内）；\tagC\ （饱和） &
非积分胶子分布过程依赖 \cite{Dominguez:2011wm}；无受控幂次修正 \\
\bottomrule
\end{tabularx}
\end{table}

\begin{table}[p]
\centering\footnotesize
\renewcommand{\arraystretch}{1.18}
\caption{QCD 因子化相关结论的严格性总表（四）：高扭度、方法论与其他（条目 21--30）。}
\label{tab:rigor-master-4}
\begin{tabularx}{\textwidth}{@{}L{0.45cm}L{2.0cm}YL{2.6cm}L{2.35cm}Y@{}}
\toprule
\# & 主题 & 已确立的结论（要点） & 关键文献 & 标签 & 主要缺口／反例 \\
\midrule
21 & 高扭度与扭度 3 &
DIS 扭度 4（OPE）；强子 $1/Q^2$ 广义因子化；扭度 3 单自旋不对称 &
\cite{Ellis:1982cd,Jaffe:1983hp,Qiu:1990xy,Qiu:1991pp,Efremov:1984ip} &
\tagR\ （DIS 矩）；\tagA\ （$1/Q^2$）；\tagP\ （扭度 3）；\tagV\ （强子 $1/Q^4$） &
与首阶扭度的重整子混合；扭度 3 无全阶证明 \\
22 & 重整子 &
IR 重整子对应幂次修正；微扰级数阶乘发散 &
\cite{tHooft:1977xjm,Mueller:1984vh,Beneke:1998ui,Bali:2014fea} &
\tagA；\tagC\ （非微扰 OPE） &
QCD 中无非微扰的 OPE 定义 \\
23 & 区域方法 &
欧氏极限展开严格；闵氏阈值、在壳、Glauber 展开仅在特例中证明 &
\cite{Smirnov:1990rz,Beneke:1997zp,Jantzen:2011nz,Semenova:2018cwy,Gardi:2022khw} &
\tagR\ （欧氏）；\tagP\ （闵氏） &
一般情形为猜想 \cite{Semenova:2018cwy} \\
24 & LaMET 与 pseudo-PDF &
准 PDF 与 pseudo-ITD 对光锥 PDF 的匹配；Wilson 线夸克双线性算符全阶乘法可重整 &
\cite{Ji:2013dva,Izubuchi:2018srq,Radyushkin:2017cyf,Ishikawa:2017faj,Ji:2017oey} &
\tagA &
依赖 QCD 类空 OPE；幂次修正 $\sim\Lambda^2/(x^2(1-x)P_z^2)$ \cite{Braun:2018brg}；批评 \cite{Rossi:2018zkn} \\
25 & 阈值重求和与再因子化 &
首阶幂次阈值因子化与 Mellin 空间指数化；次首阶幂次 LL &
\cite{Sterman:1986aj,Catani:1989ne,Becher:2006mr,Beneke:2018gvs} &
\tagA\ （LP）；\tagP\ （NLP） &
NLP 端点发散需逐例处理 \cite{Beneke:2022obx} \\
26 & 非全局对数 &
$O(\alpha_s^2)$ 起的 NGL；BMS 方程（大 $N_c$）；有限 $N_c$；EFT 因子化 &
\cite{Dasgupta:2001sh,Banfi:2002hw,Weigert:2003mm,Becher:2016mmh} &
\tagR\ （存在性）；\tagA\ （大 $N_c$ LL）；\tagP &
强子对撞中 Glauber 相位导致超领头对数（条目 13） \\
27 & DIS 重夸克方案 &
保留重夸克质量的因子化，余项一致为 $O(\Lambda/Q)$；ACOT 等方案 &
\cite{Collins:1978wz,Aivazis:1993pi,Collins:1998rz,Kramer:2000hn,Forte:2010ta} &
\tagA &
强子--强子中保留入射夸克质量不可取（Collins 2011 §14.6） \\
28 & 光生与分解光子 &
光子结构函数（点状部分）；直接光生类 DIS；分解光生类强子--强子；质子光子 PDF &
\cite{Witten:1977ju,Bardeen:1978hg,Klasen:2002xb,Manohar:2016nzj} &
\tagA\ （直接）；\tagC\ （分解）；\tagV\ （衍射分解） &
直接/分解在 LO 以上的划分不唯一；衍射光生压低 \cite{H1:2007jtx,Klasen:2008ah} \\
29 & 核 PDF 与核中因子化 &
$eA$ 同 DIS；$pA$ DY 同 DY 加 $A^{1/3}/Q^2$ 核增强扭度 4 &
\cite{Qiu:1990xy,Luo:1994np,Bodwin:1988fs,Eskola:2021nhw} &
\tagA &
要求能量远大于随靶长度增长的标度 \cite{Bodwin:1988fs}；小 $x$ 饱和区 \tagC \\
30 & 其他要点 &
$e^+e^-\to hX$ 碎裂；Sudakov 形状因子；多腿振幅红外结构；振幅层次硬--软--共线因子化 &
\cite{Collins:1981uw,Mueller:1979ih,Collins:1980ih,Becher:2009cu,Almelid:2015jia,Feige:2014wja} &
\tagA &
偶极公式作为全阶陈述在三圈失效（\tagV）；Feige--Schwartz 排除初末态共线构型 \\
\bottomrule
\end{tabularx}
\end{table}
